\documentclass{ximera}

\input{../../preamble.tex}

\title{Algebra}

\begin{document}

\begin{abstract}
positive and negative
\end{abstract}
\maketitle



Algebra's strength is in identifying zeros. 


This is largely due to the Zero Property Property, which says that 



\[
\text{ If } \, a \cdot b = 0 \, \text{ then either } a = 0 \, \text{ or } \, b = 0
\]


The first step in identifying a funciton's zeros is to form an equation by setting the function equal $0$.

Now find solutions to the equation. 


Unless you know of some handy property of the function, our procedure often begin with ``get eveything on one side and $0$ on the other and factor''.




Contrast that to increasing and decreasing, which do not involve equations to solve.  They are comparisons of movement or change between the range and domain. Our algebra is not that good at such comparisons. 


The derivative rephrases this comparison of change back into algebra, where we have methods. 





\begin{idea} \textbf{\textcolor{blue!55!black}{Movement to Arithmetic}} 

Let $f$ be a function. 

Then the derivative is denoted as $f'$. 

\begin{itemize}
\item where $f'$ is positive, $f$ is increasing  
\item where $f'$ is negative, $f$ is decreasing   
\end{itemize}

\end{idea}

This allows us to bring our algebraic tools to the question of function behavior. 










\textbf{Quadratics}



Any quadratic can be written in the form $A \, x^2 + B \, x + C$. 

The derivative is given by $2A \, x + B$.


The derivative switches signs at $\frac{-B}{2A}$, which is the only critical. It corresponds to the first coordinate of the vertex. 



\textbf{$\blacktriangleright$} If $A < 0$, then 

\begin{itemize}
\item $f' > 0$ on $\left( -\infty, \frac{-B}{2A} \right)$ and $f$ is increasing on $\left( -\infty, \frac{-B}{2A} \right)$. \\
\item $f' < 0$ on $\left( \frac{-B}{2A}, -\infty \right)$ and $f$ is decreasing on $\left( \frac{-B}{2A}, \infty \right)$. 
\end{itemize}







\textbf{$\blacktriangleright$} If $A > 0$, then 

\begin{itemize}
\item $f' < 0$ on $\left( -\infty, \frac{-B}{2A} \right)$ and $f$ is decreasing on $\left( -\infty, \frac{-B}{2A} \right)$. \\
\item $f' > 0$ on $\left( \frac{-B}{2A}, -\infty \right)$ and $f$ is increasing on $\left( \frac{-B}{2A}, \infty \right)$. 
\end{itemize}








\begin{example} Quadratics


Where does $Q(k) = 3 \, k^2 + 4 \, k -5$ increase and decrease?

\textbf{\textcolor{red!75!green}{explanation}} 


We can get the derivative of $Q$: $Q'(k) = 6 \, k + 4$, which is a linear function with a postive leading coeffcient.

The zero of $Q'$ is $-\frac{4}{6} = -\frac{2}{3}$, which makes $-\frac{2}{3}$ a critical number for $Q$.

\begin{itemize}
\item $Q'(k)$ is negative on $\left(-\infty, \frac{2}{3} \right)$. 
\item $Q'(k)$ is positive on $\left(\frac{2}{3}, \infty \right)$. 
\end{itemize}

This gives us the behavior of $Q$.

\begin{itemize}
\item $Q'(k)$ is decreasing on $\left(-\infty, \frac{2}{3} \right)$. 
\item $Q'(k)$ is increasing on $\left(\frac{2}{3}, \infty \right)$. 
\end{itemize}

This reveals a minimum value for $Q$ at $\left(\frac{2}{3}$.

Since $Q'$ switches signs over $-\frac{2}{3}$, $Q$ switches behavior. $Q$ switches from decreasing to increasing, which means $-\frac{2}{3}$ is the location of a local minimum.  For a quadratic, this is also the global minimum value of $Q$.


\end{example}

























\begin{center}
\textbf{\textcolor{green!50!black}{ooooo-=-=-=-ooOoo-=-=-=-ooooo}} \\

more examples can be found by following this link\\ \link[More Examples of the Derivative]{https://ximera.osu.edu/csccmathematics/precalculus2/precalculus2/theDerivative/examples/exampleList}

\end{center}









\end{document}
